\documentclass[11pt]{article}

\usepackage[margin=1in]{geometry}
\usepackage{amsmath,amssymb}
\usepackage{booktabs}
\usepackage{enumitem}
\usepackage{xcolor}
\usepackage{listings}
\usepackage{microtype}
\usepackage[hidelinks]{hyperref}

\definecolor{codebg}{gray}{0.96}
\lstset{
  basicstyle=\ttfamily\small,
  backgroundcolor=\color{codebg},
  breaklines=true,
  columns=fullflexible,
  keepspaces=true,
  frame=single,
  framesep=4pt,
  xleftmargin=4pt,
  literate={~}{{$\sim$}}1
}
\newcommand{\code}[1]{\texttt{#1}}
\newcommand{\rmse}{\mathrm{RMSE}}
\newcommand{\Truth}{\mathrm{Truth}}
\newcommand{\tr}{\operatorname{tr}}
\newcommand{\Var}{\operatorname{Var}}
\newcommand{\E}{\mathbb{E}}

\title{\textbf{LPS Correctness and Validation: A Precise Experimental Plan}\\[2pt]
\large Turning the pre-submission audit roadmap into implementable, confound-controlled experiments}
\author{Pre-submission audit (companion to \texttt{lps\_ps\_lps\_jasa\_style\_audit\_2026-06-09.md})}
\date{2026-06-09}

\begin{document}
\maketitle

\begin{abstract}
This document specifies, for every item in the LPS roadmap (Tiers 0--4), an experiment
or test that is detailed enough that two independent implementers would build the
\emph{same} experiment and reach the \emph{same} verdict. Each entry states a falsifiable
claim, the exact data-generating process (DGP), the exact \code{fit.lps} configuration,
the statistic computed, a numeric acceptance criterion, and the validity safeguards that
prevent the experiment from answering a confounded question. The shared conventions in
Section~\ref{sec:conventions} (RNG discipline, tolerances, the canonical DGP library, the
linear-smoother extraction procedure and its leave-one-out caveat, the paired-comparison
protocol) are defined once and referenced throughout; they are mandatory, not advisory.
The priority order is correctness-first: Tier~0 establishes that the estimator is what it
claims to be before any comparative or accuracy experiment is run.
\end{abstract}

\tableofcontents

\section{Scope and how to read this document}

The object under test is the exported estimator \code{fit.lps} (and \code{predict.lps})
in package \code{geosmooth}, with the signature and defaults
\begin{lstlisting}
fit.lps(X, y, foldid = NULL,
        support.grid = c(10L, 15L, 20L), degree.grid = 0:2,
        kernel.grid = c("gaussian", "tricube"),
        cv.folds = 5L, cv.seed = 1L, X.eval = NULL,
        coordinate.method = c("coordinates", "local.pca"),
        chart.dim = NULL,
        local.chart.method = c("pca", "second.order.svd"),
        auto.chart.support.metric = c("coordinates", "operator", "both"),
        auto.chart.selection.metric = c("coordinates", "operator"),
        backend = c("auto", "R", "cpp", "cpp.local.pca"),
        design.basis = c("orthogonal.polynomial.drop", "monomial",
                         "weighted.qr", "weighted.qr.drop"),
        design.drop.tol = 1e-8,
        ridge.multiplier.grid = c(0, 1e-10, 1e-8),
        ridge.condition.max = 1e12,
        unstable.action = c("na", "mean"),
        outcome.family = c("gaussian", "bernoulli", "binomial"))
\end{lstlisting}

PS-LPS (\code{fit.ps.lps}) is deliberately out of scope here; this plan concerns LPS
alone, per the stated priority of LPS algorithmic and implementation correctness.

Every experiment uses the same eight-field template: \textbf{Claim} (falsifiable);
\textbf{Failure mode addressed}; \textbf{DGP}; \textbf{Configuration} (every non-default
argument pinned); \textbf{Statistic}; \textbf{Acceptance criterion} (numeric);
\textbf{Validity safeguards}; \textbf{Artifacts}. ``Tests'' (Tier~0, and the regression
guards) are deterministic \code{testthat} assertions that must pass in CI. ``Studies''
(the comparative items in Tiers~1--4) are scripted experiments with predeclared decision
rules; their outputs are reports plus a machine-readable verdict row.

\section{Global conventions and shared infrastructure}
\label{sec:conventions}

\subsection{Reproducibility and RNG discipline}
\label{sec:rng}

\begin{enumerate}[leftmargin=1.4em]
\item Every stochastic quantity is generated immediately after an explicit
\code{set.seed(s)}. Seeds are recorded in the result artifact. For replicated studies the
replicate $r\in\{1,\dots,R\}$ uses seed $s_r = s_0 + r$ with a documented base $s_0$.
\item Cross-validation fold assignment is \emph{always} supplied explicitly via the
\code{foldid} argument; tests must not rely on the internal \code{cv.seed} default, because
the internal fold construction is an implementation detail that may change. For paired
comparisons the \emph{same} \code{foldid} vector is passed to both arms.
\item Each result artifact stores \code{sessionInfo()}, the \code{geosmooth} package
version and git commit hash, the BLAS/LAPACK identification, and the full argument list of
every \code{fit.lps} call. A study whose artifact lacks these is invalid and must be re-run.
\item Determinism requirement: re-running any Tier~0 test twice in the same environment
must yield bitwise-identical numeric results. Any test that is not deterministic is a
defect in the test, not a tolerance to be widened.
\end{enumerate}

\subsection{Tolerance conventions}
\label{sec:tol}

Let $a$ be a computed value and $b$ a reference. Define absolute error $|a-b|$ and relative
error $|a-b|/(|b|+\delta)$ with floor $\delta=10^{-12}$.

\begin{itemize}[leftmargin=1.4em]
\item \textbf{Algebraic identities} (e.g.\ $\hat y = Sy$, orthogonal vs.\ monomial span):
absolute tolerance $\tau_{\mathrm{alg}}=10^{-10}$.
\item \textbf{Polynomial reproduction / solver exactness:} $10^{-8}$ for
\code{design.basis} $\in\{$\code{orthogonal.polynomial.drop}, \code{weighted.qr},
\code{weighted.qr.drop}$\}$; relaxed to $10^{-6}$ for \code{monomial} on raw coordinates,
where conditioning is worse \emph{by design}. The relaxation must be justified in the test
comment by reporting the realized design condition number.
\item \textbf{Iterative solves} (local logistic IRLS): convergence tolerance $10^{-7}$ on
the coefficient step (matching the implementation), comparison tolerance $10^{-6}$.
\item \textbf{Monte-Carlo quantities} (coverage, calibration, rates): no fixed tolerance;
the acceptance criterion is a confidence interval or hypothesis test with the replicate
count $R$ chosen so the Monte-Carlo standard error is below one third of the decision
threshold (state the power calculation in the script header).
\end{itemize}

\subsection{Canonical data-generating processes (DGP library)}
\label{sec:dgp}

These are defined once. Every experiment references a DGP by tag and overrides only the
parameters it names. All coordinates are double precision; all latent draws use
Section~\ref{sec:rng}.

\begin{description}[leftmargin=0em,style=nextline]
\item[G1 — ambient polynomial.] Ambient $D\in\{2,3\}$. $X_i \overset{iid}{\sim}
\mathrm{Unif}([-1,1]^D)$, $i=1,\dots,n$. Truth a fixed polynomial of degree $p$; for
$D=2,p=2$ use $f(x)=0.5+1.0x_1-0.7x_2+0.4x_1^2+0.3x_1x_2-0.6x_2^2$. Noiseless unless stated.
Used for ambient reproduction and df checks.
\item[G2 — flat embedded subspace.] Intrinsic $d$, ambient $D>d$. Latent
$u_i\sim\mathrm{Unif}([-1,1]^d)$; a fixed random orthonormal $Q\in\mathbb{R}^{D\times d}$
(seed recorded); $X_i = Qu_i$. Truth a degree-$p$ polynomial in the \emph{intrinsic} $u$.
The embedding is a linear isometry, so a degree-$p$ local fit in the tangent chart
reproduces $f$ exactly. Used for intrinsic reproduction.
\item[G3a — paraboloid (known curvature).] $d=2,D=3$. $u\sim\mathrm{Unif}$ on the disk of
radius $\rho_0=1$; $X=(u_1,u_2,(u_1^2+u_2^2)/(2R))$. Principal curvature $\kappa=1/R$ at the
apex; $R$ is the curvature knob. Truth options: $f_{\mathrm{lin}}(u)=u_1$ (intrinsic-linear,
isolates curvature bias) and $f_{\mathrm{smooth}}(u)=\sin(\pi u_1)\cos(\pi u_2)$.
\item[G3b — sphere cap (constant curvature).] $d=2,D=3$. A cap of a sphere of radius $R$:
latent $u\sim\mathrm{Unif}$ on the planar disk $\lVert u\rVert\le\rho_0<R$, with
$X=(u_1,u_2,\,R-\sqrt{R^2-\lVert u\rVert^2})$. Constant Gaussian curvature $1/R^2$ --- a
constant-positive-curvature complement to G3a's position-varying paraboloid curvature. Truth
options as in G3a.
\item[G3c — 1-D curve.] $d=1,D=3$. Helix $X(t)=(\cos t,\sin t, c\,t)$, $t\sim\mathrm{Unif}[0,2\pi)$.
\item[G3d — torus patch.] $d=2,D=3$, standard torus with radii $(R_1,R_2)$, latent angles on
a sub-rectangle to avoid wrap-around.
\item[G4 — stratified / varying dimension.] Two glued strata sharing a boundary: stratum A is
a 1-D segment $\{(t,0,0):t\in[-1,0]\}$ (thickened by noise $\eta$), stratum B is a 2-D patch
$\{(s_1,s_2,0):s_1\in[0,1],s_2\in[-0.5,0.5]\}$. True local intrinsic dimension is $1$ on A and
$2$ on B. Used to verify a dimension stabilizer does \emph{not} blur the true boundary.
\item[G5 — clustered / repeated measures.] $K$ clusters; centers $c_k\sim\mathrm{Unif}([-1,1]^2)$;
within cluster, $m$ points $x=c_k+\zeta$, $\zeta\sim N(0,\sigma_x^2 I)$. Response
$y=g(x)+b_k+\varepsilon$, $b_k\sim N(0,\tau^2)$, $\varepsilon\sim N(0,\sigma^2)$, intra-class
correlation $\rho=\tau^2/(\tau^2+\sigma^2)$. Used for grouped-CV validity.
\item[G6 — binary surface.] Smooth bounded log-odds $\eta(x)$ with offset $\alpha$ chosen so
$\E[p]$ hits a target prevalence; $p(x)=\mathrm{expit}(\alpha+\eta(x))$ clipped by construction
to $[0.05,0.95]$ (no spike profiles); $Y_i\sim\mathrm{Bernoulli}(p(x_i))$. Default
$\eta(x)=1.5\sin(\pi x_1)$ on $x\in[-1,1]^2$.
\item[G7 — compositional / structural zeros.] $X_i$ a $D$-part composition: draw
$\mathrm{Dirichlet}(\boldsymbol\alpha)$, set a documented subset of parts to exact $0$, renormalize
to the simplex. Used to characterize behavior on 16S-like inputs.
\end{description}

\subsection{Linear-smoother extraction, degrees of freedom, and the leave-one-out caveat}
\label{sec:smoother}

Several experiments rely on the fact that, \emph{for a single fixed configuration} $\theta$
(one support size, one degree, one kernel, a fixed \code{chart.dim}, a fixed ridge), the LPS
fitted vector is linear in the response:
\begin{equation}
\hat y = S(\theta)\,y, \qquad S(\theta)\ \text{depends only on } X,\ \text{weights},\ \theta,\ \text{not on } y.
\label{eq:linear}
\end{equation}

\paragraph{Mandatory restriction.} Equation~\eqref{eq:linear} holds only when no
data-driven selection occurs. Therefore every experiment that uses $S$ must call
\code{fit.lps} with \emph{singleton} grids (\code{support.grid}, \code{degree.grid},
\code{kernel.grid} of length one) and an explicit numeric \code{chart.dim} (never
\code{"auto"}/\code{"local.auto"}, whose estimate, although $y$-independent, is a separate
map). Using the CV-selected pipeline and then asserting linearity is a known trap: the
selected $\hat\theta(y)$ depends on $y$, so the end-to-end pipeline is \emph{not} linear,
and the assertion would (correctly) fail for the wrong reason.

\paragraph{Extraction procedure.} With $X.\mathrm{eval}=X$, fixed $\theta$, and the chosen
\code{design.basis}: column $j$ of $S$ is $S_{\cdot j} = \texttt{fit}(e_j)$ where $e_j$ is the
$j$-th unit response vector (linearity makes finite differencing exact, so no step size is
needed). Assemble $S\in\mathbb{R}^{n\times n}$, then $\hat y = Sy$, effective degrees of
freedom $\mathrm{df}=\tr S$, and pointwise sensitivities $S_{ii}$.

\paragraph{Two distinct leave-one-out notions (do not conflate).}
\begin{enumerate}[leftmargin=1.4em]
\item \emph{Response-removal LOO}: hold the geometry fixed, set $y_i$'s influence to zero.
For a linear smoother this is \emph{exactly} $\hat y^{(-i)}_{\mathrm{resp}}(x_i)=(\hat y_i - S_{ii}y_i)/(1-S_{ii})$.
\item \emph{Training-removal LOO}: delete observation $i$ from the training set and refit.
For a $k$-NN-bandwidth smoother this \emph{changes the support} of $x_i$ (its $k$-th neighbor
and hence the bandwidth $h=\max$ distance change), so it is \emph{not} equal to the
response-removal shortcut.
\end{enumerate}
The honest generalization target is training-removal LOO; the cheap shortcut is
response-removal. Their gap is a real property of $k$-NN local regression and must be
\emph{measured}, not assumed away (see E0.3b). Asserting the $(1-S_{ii})$ identity against a
training-removal refit is invalid and is explicitly prohibited.

\subsection{Paired-comparison protocol}
\label{sec:paired}

Any ``method A vs.\ method B'' accuracy comparison must hold fixed, across the two arms:
the geometry $X$, the truth $f$ (or $p$), the realized response $y$ (same noise draw), the
\code{foldid}, and—unless the comparison is explicitly \emph{about} that axis—the support
size, degree, kernel, and \code{chart.dim}. The unit of analysis is the per-case delta
$\Delta = M_A - M_B$ of the target metric. Report the sign-test / Wilcoxon signed-rank
across cases and replicates, the median $\Delta$ with a bootstrap CI, and the full per-case
table. A study that lets both arms re-select tuning by CV (so the arms differ on multiple
axes at once) may be run only \emph{in addition} to the matched comparison, never as its
replacement.

\subsection{Decision-rule conventions}
\label{sec:decision}

Each study predeclares: the primary metric, the effect-size threshold that constitutes
``materially better'', the replicate count $R$ with its Monte-Carlo-error justification, and
the statistical test. ``Promote'' / ``do not promote'' verdicts are emitted as a one-row
machine-readable summary. Post-hoc threshold changes invalidate the study.

\section{Tier 0 — Statistical-correctness battery (deterministic tests)}

These are the cheapest possible proofs that the estimator is what it claims to be. They are
\code{testthat} tests, must be deterministic (Section~\ref{sec:rng}), and gate everything
downstream. None of them involves model selection.

\subsection{E0.1 — Polynomial reproduction (ambient and intrinsic)}

\paragraph{Claim.} A degree-$p$ LPS reproduces any global polynomial of degree $\le p$ with
error at machine-solve precision, for every kernel and every \code{design.basis}, and—on a
\emph{flat} manifold—for the \code{local.pca} chart at the correct intrinsic dimension.

\paragraph{Failure mode addressed.} Centering, design-matrix construction, or
chart-coordinate bugs that introduce bias even on data the model can represent exactly. This
is the single most diagnostic LPS test and is currently absent.

\paragraph{DGP.} (a) Ambient: G1 with $D\in\{2,3\}$, $p\in\{1,2\}$, $n=200$, noiseless.
(b) Intrinsic: G2 with $d\in\{1,2\}$, $D=d+2$, $p\in\{1,2\}$, $n=200$, noiseless.

\paragraph{Configuration.} Singleton grids; support large enough for full rank.
\begin{lstlisting}
# (a) ambient reproduction
fit.lps(X, y, foldid = rep(1:2, length.out = n),
        support.grid = K, degree.grid = p, kernel.grid = ker,
        coordinate.method = "coordinates", backend = "R",
        design.basis = basis, ridge.multiplier.grid = 0,
        ridge.condition.max = Inf, unstable.action = "na")
# (b) intrinsic reproduction: coordinate.method = "local.pca", chart.dim = d
\end{lstlisting}
Sweep \code{ker} over all four kernels, \code{basis} over all four design bases, with
$K = 3\cdot c_p$ where $c_p$ is the design-column count $c_p=1+d+\binom{d+1}{2}\mathbf{1}\{p=2\}$
(ambient uses $d=D$). The factor $3$ guarantees the support is over-determined for the degree.

\paragraph{Statistic.} $\max_i |\hat f(x_i)-f(x_i)|$ over $X.\mathrm{eval}=X$.

\paragraph{Acceptance criterion.} $<10^{-8}$ for orthogonal/QR bases, $<10^{-6}$ for
\code{monomial} (Section~\ref{sec:tol}), for \emph{all} kernel $\times$ basis $\times$ degree
combinations. The test also asserts the reported minimum design rank equals $c_p$.

\paragraph{Validity safeguards.} (i) The intrinsic case uses a \emph{flat} (isometric)
embedding only: on a curved manifold exact reproduction is false (curvature bias), so using
G3 here would be a spurious failure. (ii) $X$ drawn from a continuous distribution with a
fixed seed so supports are in general position (no accidental collinearity); the rank
assertion guards this. (iii) \code{chart.dim} set to the true $d$, never \code{"auto"}, so
the test isolates reproduction from dimension estimation.

\paragraph{Artifacts.} \code{testthat} file \code{test-lps-reproduction.R}; on failure,
dump the offending neighborhood's design, weights, and condition number.

\subsection{E0.2 — Linear-smoother identity and degrees of freedom}

\paragraph{Claim.} For a fixed configuration, (i) $\hat y = Sy$ exactly; (ii) $S$ is
independent of $y$; (iii) $\mathrm{df}=\tr S$ equals the finite-difference trace
$\sum_i \partial \hat y_i/\partial y_i$.

\paragraph{Failure mode addressed.} Any hidden $y$-dependence in the weighting or support
(which would break every downstream df / variance / CV-shortcut claim), and an incorrect df
accounting.

\paragraph{DGP.} G1, $D=2$, $n=120$, $y$ realized once with $N(0,1)$ entries (the values are
irrelevant; only linearity is tested).

\paragraph{Configuration.} Singleton grids; \code{coordinate.method} both
\code{"coordinates"} and \code{"local.pca"} with explicit \code{chart.dim=2};
\code{design.basis = "orthogonal.polynomial.drop"}; fixed ridge.

\paragraph{Statistic.} Build $S$ per Section~\ref{sec:smoother}. Compute
(i) $\max|{\texttt{fit}(y)}-Sy|$ for three independent random $y$; (ii) $\max|S^{(y_a)}-S^{(y_b)}|$
for two different base responses; (iii) $|\tr S - \sum_i (\texttt{fit}(y+ \epsilon e_i)_i-\texttt{fit}(y)_i)/\epsilon|$
with $\epsilon=1$ (exact by linearity).

\paragraph{Acceptance criterion.} (i) $<\tau_{\mathrm{alg}}=10^{-10}$; (ii) $<10^{-12}$
(literally the same matrix); (iii) $<10^{-10}$.

\paragraph{Validity safeguards.} Singleton grids and numeric \code{chart.dim} per
Section~\ref{sec:smoother}; assert that the run performed no candidate selection (CV table
has one row).

\paragraph{Artifacts.} \code{test-lps-linear-smoother.R}; persist one $S$ matrix as a fixture
for E0.3 and E4.

\subsection{E0.3a — Response-removal LOO shortcut (algebraic)}

\paragraph{Claim.} Whatever leave-one-out / GCV residual the LPS (or GCV) code reports matches
the analytic shortcut $(y_i-\hat y_i)/(1-S_{ii})$ formed from the \emph{independently} extracted
$S$ of E0.2. This is a unit test of the \textbf{GCV/LOO formula implementation}, not a validation
of the estimator: response-removal LOO equals that shortcut \emph{by derivation}, so it cannot
test anything beyond whether the code computes the formula correctly. The scientifically
meaningful leave-one-out question --- training-removal, where the $k$-NN support changes --- is
E0.3b, which carries the weight.

\paragraph{Failure mode addressed.} An incorrect GCV/LOO algebra in the code that would silently
bias any $S$-based selection.

\paragraph{DGP / Configuration.} As E0.2.

\paragraph{Statistic.} Compare the code-reported per-point LOO/GCV residual to
$(y_i-\hat y_i)/(1-S_{ii})$ with $S$ from E0.2. If the package has no separate GCV/LOO path for
LPS, this test has no independent target: it is then deferred (folded into E0.2) and must
\emph{not} be reconstructed purely from $S$ and then ``verified'' against itself --- that would
be tautological.

\paragraph{Acceptance criterion.} $\max_i$ relative error $<10^{-8}$ against the independent code
path (or deferred if no such path exists); rows with $S_{ii}\ge 1-10^{-6}$ are reported and
excluded (near-interpolation).

\paragraph{Validity safeguards.} Do \emph{not} compare against a training-removal refit here
(Section~\ref{sec:smoother}).

\subsection{E0.3b — Training-removal vs.\ shortcut gap (characterization study)}

\paragraph{Claim (characterization, not pass/fail).} Quantify how far the cheap response-removal
shortcut is from honest training-removal LOO for $k$-NN-bandwidth LPS, as a function of $k$.

\paragraph{Failure mode addressed.} Silently using the $(1-S_{ii})$ shortcut as if it were
honest CV when the $k$-NN support change makes it optimistic.

\paragraph{DGP.} G3a ($f_{\mathrm{smooth}}$), $n=400$, noise $\sigma=0.05$; $k\in\{8,15,30,60\}$;
$R=20$ replicates.

\paragraph{Configuration.} Fixed degree $1$, \code{chart.dim=2}, \code{gaussian};
training-removal LOO for a random subsample of $200$ indices via re-fitting on $X[-i]$ and
predicting $x_i$.

\paragraph{Statistic.} Per $i$: $g_i=\hat y^{(-i)}_{\mathrm{train}}(x_i)-\hat y^{(-i)}_{\mathrm{resp}}(x_i)$;
report median and 90th percentile of $|g_i|/\mathrm{sd}(y)$ by $k$.

\paragraph{Acceptance criterion (advisory threshold).} If median relative gap $<0.05$ for the
$k$ used in practice, the shortcut may stand in for honest CV in reports; otherwise reports
must use training-removal CV (or grouped CV, E1.10). The verdict is recorded, not asserted.

\paragraph{Validity safeguards.} Same noise draw across the two LOO notions; exclude
near-interpolation rows; hold $k$ fixed within a cell.

\subsection{E0.4 — Boundary-bias property (local linear vs.\ local constant)}

\paragraph{Claim.} Local linear (degree $1$) is first-order boundary-bias-free whereas local
constant (degree $0$) carries $O(h)$ boundary bias: near the domain edge the degree-$1$ error
is a small multiple of its interior error, while the degree-$0$ boundary/interior error ratio
is large.

\paragraph{Failure mode addressed.} A degree-$1$ implementation that has lost the
boundary-correction property (e.g.\ through mis-centering).

\paragraph{DGP.} 1-D design $x_i$ on $[0,1]$, $n=300$ (equally spaced plus a fixed jitter,
seed recorded); truth $f(x)=e^{x}$ (nonzero slope and curvature at the boundary); noiseless.

\paragraph{Configuration.} \code{coordinate.method="coordinates"}, singleton $K=30$,
\code{gaussian}, degrees $0$ and $1$ in separate fits, \code{ridge.multiplier.grid=0}.

\paragraph{Statistic.} Boundary set $B=\{x_i: x_i< h\ \text{or}\ x_i>1-h\}$ with $h$ the
realized $K$-NN bandwidth at the edge; interior set $I$ the complement. Compute
$\bar e_B^{(p)}=\mathrm{mean}_{B}|\hat f-f|$ and $\bar e_I^{(p)}$ for $p\in\{0,1\}$.

\paragraph{Acceptance criterion.} $\bar e_B^{(1)}/\bar e_B^{(0)} < 0.5$ \emph{and}
$\bar e_B^{(1)}/\bar e_I^{(1)} < 3$ \emph{and} $\bar e_B^{(0)}/\bar e_I^{(0)} > 5$.

\paragraph{Validity safeguards.} Noiseless (isolates bias); same $x$, $h$, kernel across the
degree arms (paired); truth chosen nonlinear so degree~1 does not reproduce it exactly (which
would make all errors $0$ and the ratio undefined).

\subsection{E0.5 — Consistency and rate (mean-response)}

\paragraph{Claim.} With the bandwidth schedule below, Truth-RMSE decreases at the local-linear
manifold rate $\rmse \asymp n^{-2/(d+4)}$ on a curved $d$-manifold.

\paragraph{Failure mode addressed.} A bias floor masquerading as inconsistency (the trap of
fixing $k$ so bias never vanishes), or an estimator that fails to converge.

\paragraph{DGP.} G3a ($f_{\mathrm{smooth}}$, $d=2$), noise $\sigma=0.1$;
$n\in\{200,400,800,1600,3200\}$; $R=30$ replicates per $n$.

\paragraph{Configuration.} Fixed degree $1$, \code{chart.dim=2}, \code{gaussian},
\code{design.basis="orthogonal.polynomial.drop"}. Support grows as
$k(n)=\mathrm{round}\!\big(c\,n^{4/(d+4)}\big)$ (for $d=2$: $k\propto n^{2/3}$), with $c$ fixed
so $k(200)=20$; \code{support.grid=k(n)} (singleton, so no selection).

\paragraph{Statistic.} $\rmse_p(n)=\big(\tfrac1n\sum_i(\hat f(x_i)-f(x_i))^2\big)^{1/2}$,
averaged over replicates; fit $\log \rmse = a + b\log n$.

\paragraph{Acceptance criterion.} (Smoke, required) the 95\% CI for $b$ lies entirely below
$-0.1$. (Rate, recommended) the CI for $b$ contains the theoretical $-2/(d+4)=-1/3$, or the
deviation is explained by a stated finite-sample/boundary argument.

\paragraph{Validity safeguards.} Bandwidth schedule pinned to the optimal growth (otherwise a
bias floor invalidates the rate); same seed structure across $n$; report bias and variance
components separately so a plateau is correctly attributed.

\subsection{E0.6 — Binary recovery and calibration}

\paragraph{Claim.} Both binary modes produce probability estimates that (i) are consistent
($\rmse_p\!\downarrow$ with $n$) and (ii) are calibrated at large $n$ (calibration slope
$\approx 1$, intercept $\approx 0$).

\paragraph{Failure mode addressed.} Mis-scaled or mis-clipped probabilities; a selection
metric that does not target probability quality; silent degradation under class imbalance.

\paragraph{DGP.} G6 on $x\in[-1,1]^2$ (treated as \code{coordinates}), prevalence
$\E[p]\in\{0.1,0.3,0.5\}$, $n\in\{500,1000,2000,4000\}$, $R=40$.

\paragraph{Configuration.} \code{outcome.family} $\in\{$\code{"bernoulli"},\code{"binomial"}$\}$,
\code{support.grid}/\code{degree.grid}/\code{kernel.grid} as the routine grids (selection
\emph{is} allowed here—this is an accuracy study, not an $S$ study), \code{backend="R"}.

\paragraph{Statistic.} $\rmse_p$ vs.\ the known $p$; reliability via $10$ equal-count bins of
$\hat p$ (observed frequency vs.\ mean $\hat p$); calibration slope/intercept from the logistic
regression $Y\sim \mathrm{logit}(\hat p)$ on a held-out half.

\paragraph{Acceptance criterion.} (i) the $\log\rmse_p$ vs.\ $\log n$ slope CI lies below
$-0.1$ at every prevalence; (ii) at $n=4000$, calibration slope $\in[0.8,1.25]$ and intercept
$\in[-0.25,0.25]$. The Bernoulli-vs-binomial difference is \emph{reported}, not asserted.

\paragraph{Validity safeguards.} $p$ bounded in $[0.05,0.95]$ (no spike profiles); calibration
fit on data disjoint from the LPS fit; at prevalence $0.1$ with $n=500$ expect logistic
fallbacks—stratify all metrics by \code{logistic.diagnostics} fallback fraction and report it.

\subsection{E0.7 — Cross-validation has no response leakage}

\paragraph{Claim.} Perturbing a held-out response $y_i$ leaves $i$'s own out-of-fold
prediction exactly unchanged.

\paragraph{Failure mode addressed.} Any leakage of held-out responses into their own
out-of-fold predictions (which would make CV optimistic).

\paragraph{DGP.} G1, $D=2$, $n=100$, explicit \code{foldid} with $5$ folds.

\paragraph{Configuration.} Fixed configuration (singleton grids) so no selection confounds the
test. Out-of-fold prediction of fold $f$ obtained by \code{fit.lps} on
$\{j:\code{foldid}_j\ne f\}$ with $X.\mathrm{eval}=\{j:\code{foldid}_j=f\}$.

\paragraph{Statistic.} For a held-out index $i$ (in fold $f$): refit after replacing $y_i$ by
$y_i+\Delta$ ($\Delta=10$) and recompute the out-of-fold prediction at $i$. Report
$|\hat y^{\mathrm{oof}}_i(\text{perturbed})-\hat y^{\mathrm{oof}}_i(\text{original})|$.

\paragraph{Acceptance criterion.} $<10^{-12}$ (exactly zero up to roundoff). As a positive
control, perturbing $y_i$ \emph{does} change the prediction of some $j$ in another fold that has
$i$ in its training support (assert change $>0$), proving the test is not vacuous.

\paragraph{Validity safeguards.} Fixed config; the geometry (and any \code{auto} dimension) is
$y$-independent, so the test isolates response leakage.

\subsection{E0.8 — Degenerate-geometry pathologies}

\paragraph{Claim.} On each pathology, LPS terminates, returns finite predictions where defined,
and exhibits the specified guarded behavior (no silent mean fallback when
\code{unstable.action="na"}; correct diagnostics).

\paragraph{Failure mode addressed.} Crashes, silent fallbacks, or NaNs on inputs that
real omics data routinely produce.

\paragraph{DGP / cases (each a separate test).}
\begin{enumerate}[leftmargin=1.4em,nosep]
\item \emph{Duplicate points}: replicate a single $x$ value $k$ times in a support. Expect
finite prediction; a fully-degenerate support collapses to the weighted mean (degree~$0$
behavior) and the rank diagnostic equals $1$.
\item \emph{Collinear support}: points on a line in $\mathbb{R}^2$ with degree $2$. Expect
rank-deficient quadratic dropped by \code{orthogonal.polynomial.drop}; with
\code{unstable.action="na"} a configuration that cannot be solved yields \code{NA} (and is thus
avoided by selection), never a silent mean.
\item \emph{Constant response}: $y\equiv c$. Expect $\hat f\equiv c$ within $10^{-10}$ for all
degrees, kernels, bases.
\item \emph{Distance ties}: symmetric $X$ with exact ties. Expect deterministic support choice
(index tie-break) and bitwise-identical results across repeated runs.
\item \emph{Class imbalance (binary)}: G6 at prevalence $0.02$, $n=300$, \code{outcome.family="binomial"}.
Expect no crash; \code{logistic.diagnostics} reports the event-rate fallback count; fitted
probabilities $\in[0,1]$; selection either returns a finite log-loss or fails fast with the
documented ``no finite selection score'' error.
\item \emph{Compositional / structural zeros}: G7, $D=6$. Expect finite predictions; record the
fraction of supports with zero bandwidth ($h=0$) and the fallback handling.
\end{enumerate}

\paragraph{Acceptance criterion.} Per case as stated; all assertions deterministic.

\paragraph{Validity safeguards.} Test both \code{unstable.action} settings where relevant so the
guarded-vs-legacy behavior is pinned; assert on diagnostics, not only on non-crashing.

\paragraph{Artifacts.} \code{test-lps-degenerate.R}; one fixture per case.

\section{Tier 1 — Honest selection and bandwidth}

\subsection{E1.9 — Decouple bandwidth from support size}

\paragraph{Claim.} (a) The current kernel weighting ties the effective bandwidth to the
$K$-NN radius, so ``gaussian'' is near-flat (effective sample size $\approx K$) and compact
kernels zero the $K$-th neighbor (effective support $K-1$); (b) a bandwidth multiplier $b$ that
sets $h=b\cdot(\text{$K$-th NN distance})$ reproduces current behavior exactly at $b=1$;
(c) adding $b$ to the selection grid does not worsen, and on curved truths improves, Truth-RMSE.

\paragraph{Failure mode addressed.} Kernel selection silently selecting an effective-sample-size
artifact; an unexposed bandwidth that prevents reaching the curvature-relevant scale.

\paragraph{DGP.} (a)/(b) deterministic on a single fixed support (any G1 neighborhood).
(c) G3a and G3d, $n=600$, $\sigma\in\{0.03,0.1\}$, $R=30$, paired (Section~\ref{sec:paired}).

\paragraph{Configuration / statistic.}
\begin{enumerate}[leftmargin=1.4em,nosep]
\item Characterization (regression test): call the internal weight function on a fixed
distance vector for each kernel; compute Kish effective sample size
$\mathrm{ESS}=(\sum w)^2/\sum w^2$ and the $K$-th weight. Assert
$\mathrm{ESS}/K>0.9$ for \code{gaussian}, $\mathrm{ESS}/K<0.85$ for \code{tricube}, and
$w_{(K)}/\max_j w_j<10^{-6}$ for all compact kernels. This \emph{pins current behavior} so a
future change is visible.
\item Backward-compatibility (after implementing $b$): assert fits with $b=1$ equal the
current fits within $\tau_{\mathrm{alg}}$.
\item Benefit study: arms = \{select over $(K)$ only\} vs.\ \{select over $(K,b)$, $b\in\{0.5,1,2,4\}$\};
metric Truth-RMSE; predeclared ``material'' $=$ median $\Delta>5\%$ with Wilcoxon $p<0.05$.
\end{enumerate}

\paragraph{Acceptance criterion.} (a) assertions hold; (b) $b=1$ exactness; (c) verdict
recorded; promotion only if the material threshold is met and runtime cost is reported.

\paragraph{Validity safeguards.} The benefit study is matched on $(K,\text{degree},\text{kernel},
\code{chart.dim})$ except the bandwidth axis; same noise draw; the characterization uses the
\emph{actual} internal weighting routine, not a re-implementation.

\subsection{E1.10 — Nested and grouped cross-validation}

\paragraph{Claim.} (a) The selected-minimum CV score is optimistic; nested CV corrects it and
tracks held-out test error. (b) Random $K$-fold underestimates error under cluster dependence;
leave-cluster-out CV closes the gap.

\paragraph{Failure mode addressed.} Reporting a selection-optimistic number on real data; CV
leakage through dependent (subject/replicate) samples.

\paragraph{DGP.} (a) G3a, $n=800$ train $+$ $n_{\mathrm{test}}=4000$ independent test, $R=40$.
(b) G5 with $K=40$ clusters, $m=20$, $\rho\in\{0.3,0.6\}$; an independent test set of fresh
clusters; $R=40$.

\paragraph{Configuration.} Routine grids (selection is the object of study). Outer loop: $5$
folds; inner loop: $5$ folds for selection. For (b) compare \code{foldid} built by random
assignment vs.\ by cluster id.

\paragraph{Statistic.} (a) $|\,\widehat{\rmse}_{\bullet}-\rmse_{\mathrm{test}}\,|/\rmse_{\mathrm{test}}$
for $\bullet\in\{\text{selected-min},\ \text{nested}\}$. (b) the same for random-fold and
cluster-fold estimates against fresh-cluster test error.

\paragraph{Acceptance criterion.} (a) nested relative error $<0.10$ \emph{and} nested $\ge$
selected-min in expectation (optimism sign correct). (b) random-fold relative error exceeds
cluster-fold relative error by a predeclared margin ($>0.10$ at $\rho=0.6$); cluster-fold within
$0.10$ of truth.

\paragraph{Validity safeguards.} Test clusters disjoint from training clusters (no cluster
appears in both); identical truth across arms; report the realized $\rho$.

\subsection{E1.11 — Stabilizing the local intrinsic-dimension estimate}

\paragraph{Claim.} (a) Per-anchor dimension estimates from small supports are high-variance;
(b) shrinkage-to-global and graph-smoothing of the dimension field reduce misclassification
\emph{without} blurring a true dimension boundary and without worsening Truth-RMSE; (c) a
points-per-parameter gate prevents pathological degree-$2$-in-high-dimension selections.

\paragraph{Failure mode addressed.} The documented harm from dimension over-estimation
(the FB01 finding); a stabilizer that over-smooths across genuine dimension changes.

\paragraph{DGP.} (a)/(b) homogeneous: G3a ($d=2$), $n\in\{200,800\}$, $K\in\{10,20,40\}$, $R=50$.
(b) boundary safeguard: G4 (true dim $1$ on stratum A, $2$ on B). (c) any G2 with $d=3,D=5$.

\paragraph{Configuration.} \code{coordinate.method="local.pca"}, \code{chart.dim="local.auto"};
record \code{chart.dim.by.eval}. Stabilizers implemented as post-processors of the per-anchor
estimate: shrink-to-global (mode/median of the field) and $k$-NN majority vote.

\paragraph{Statistic.} (a) misclassification rate $\frac1n\sum_i\mathbf{1}\{\hat d_i\ne d\}$
and its across-replicate variance. (b) on G4, the misclassification rate \emph{within} each
stratum and at anchors whose support straddles the boundary. (c) whether the gated selector
ever returns degree $2$ with $K<c\cdot c_2$ (set $c=1.5$).

\paragraph{Acceptance criterion.} (a) report the variance decay with $K$ (characterization).
(b) a stabilizer is acceptable only if it reduces homogeneous-region misclassification by a
predeclared margin ($>30\%$ at $K=10$) \emph{and} does not increase boundary-straddling
misclassification on G4 by more than $5$ percentage points \emph{and} is non-inferior on
Truth-RMSE (CI of $\Delta$ excludes a $>2\%$ degradation). (c) with the gate, the offending
candidates receive a non-finite/penalized score; without it, they can be selected (positive
control).

\paragraph{Validity safeguards.} The G4 boundary test is mandatory: a stabilizer that looks
good only on homogeneous data but smears the stratified boundary is rejected. Dimension truth
on G4 is defined per stratum and recorded with the geometry.

\section{Tier 2 — Binary path and numerical hygiene}

\subsection{E2.12 — Binary selection metric consistency and log-loss clipping}

\paragraph{Claim.} (a) Bernoulli-mode selection currently scores \emph{unclipped} predictions
while deployment is clipped, and this can change the selected candidate; the fix is to score the
deployed (clipped) metric. (b) A log-loss truncation of $10^{-15}$ makes the selection score
unstable; $10^{-6}$ is stable.

\paragraph{Failure mode addressed.} Train/deploy metric mismatch; a high-variance selection
criterion dominated by a single confident error.

\paragraph{DGP.} (a) a constructed case: G6 with a sharp region forcing some raw predictions
outside $[0,1]$, $n=400$, where clipped vs.\ unclipped Brier rank two candidates differently.
(b) G6 with one deliberately confident-wrong held-out point.

\paragraph{Statistic.} (a) selected candidate under clipped vs.\ unclipped Brier. (b) selected
candidate and log-loss value under clip $\in\{10^{-15},10^{-6},10^{-3}\}$.

\paragraph{Acceptance criterion.} (a) a regression test asserting the selection score equals the
deployed (clipped) metric after the fix; the pre-fix discrepancy is demonstrated in the same
file as a documented motivating case. (b) the \emph{score} at clip $10^{-15}$ is demonstrably
dominated by a single confident-wrong point (the actual defect); adopt and pin a documented
finite clip ($10^{-6}$). Selection \emph{stability} across $\{10^{-6},10^{-3}\}$ is reported as a
diagnostic, \textbf{not gated} --- near-ties may legitimately reselect, so cross-clip invariance
is a robustness goal, not a hard correctness criterion.

\paragraph{Validity safeguards.} Construct the (a) case so the ranking flip is real (verify by
hand), not a tie; hold all else fixed.

\subsection{E2.13 — Ridge penalty structure alignment}

\paragraph{Claim.} In the orthogonal basis the ridge currently penalizes the constant direction,
so as the ridge grows the prediction shrinks toward $0$ rather than toward the local weighted
mean; after aligning the penalty (leave the constant direction unpenalized), large ridge shrinks
toward the local weighted mean, and tiny ridge is prediction-invariant.

\paragraph{Failure mode addressed.} A statistically wrong shrinkage target hidden behind
currently-tiny ridge multipliers.

\paragraph{DGP.} G1, $D=2$, $n=150$, mild noise.

\paragraph{Configuration.} Fixed config; \code{design.basis="orthogonal.polynomial.drop"};
sweep \code{ridge.multiplier.grid} as singletons $\rho\in\{0,10^{-8},10^{-2},1,10^2\}$ (with
\code{ridge.condition.max=Inf} so the chosen $\rho$ is actually applied).

\paragraph{Statistic.} For each anchor, $\hat f_\rho(x_i)$; compare $\hat f_{\rho\to\infty}$ to
the local weighted mean $\bar y_i^{w}$ and to $0$.

\paragraph{Acceptance criterion.} After the fix: $|\hat f_{\rho=10^2}-\bar y^{w}|$ small relative
to $|\hat f_{\rho=10^2}-0|$ (shrinks to the local mean, not to zero); and
$|\hat f_{\rho=10^{-8}}-\hat f_{\rho=0}|<10^{-6}$ (tiny-ridge invariance). A pre-fix test
documents the shrink-to-zero behavior so the change is intentional and visible.

\paragraph{Validity safeguards.} \code{ridge.condition.max=Inf} so the guard does not silently
skip the large-$\rho$ candidate; paired across $\rho$ on the same data.

\subsection{E2.14 — Local logistic robustness (separation)}

\paragraph{Claim.} The local logistic IRLS is monotone and bounded under near-separation: with
step-halving the deviance is non-increasing across iterations, coefficients and fitted
probabilities stay finite and in range, and non-convergence triggers a deterministic, telemetered
fallback.

\paragraph{Failure mode addressed.} IRLS divergence/oscillation under separation, NaNs, or a
``win'' that is actually a mass of event-rate fallbacks.

\paragraph{DGP.} A single near-separable support: $z\in\mathbb{R}$ with $y=\mathbf{1}\{z>0\}$ plus
one flipped label to avoid exact separation; weights from \code{gaussian}.

\paragraph{Configuration.} Call the local logistic solver (degree $1$, \code{outcome.family="binomial"},
\code{backend="R"}) directly on the constructed support; record per-iteration deviance.

\paragraph{Statistic.} Per-iteration deviance sequence; final coefficients; fitted probability;
telemetry (\code{converged}, \code{fallback.path}, \code{event.rate.fallback}).

\paragraph{Acceptance criterion.} With step-halving, deviance is non-increasing to within
$10^{-8}$ per step; final $\hat p\in(0,1)$ and $|\hat\beta|<\infty$; if the convergence tolerance
is not met within the iteration cap, the documented fallback fires and is recorded. Exact
separation (no flipped label) must hit the fallback, not loop or NaN.

\paragraph{Validity safeguards.} Test both the near-separable and exactly-separable cases; assert
on the deviance \emph{trajectory}, not only the endpoint, so oscillation is caught.

\section{Tier 3 — Curvature-chart re-evaluation, done correctly}

The H4/H5 ``no gains'' result for \code{second.order.svd} is confounded: both arms selected
degree~$2$ (the diagnostics' own top predictor of the effect is the cross-arm degree
difference), the error budget was noise-dominated ($\sigma\!\approx\!$ Truth-RMSE), curvature was
estimated from $\le 18$-point supports, and the metric averaged over the whole manifold. A
degree-$2$ polynomial in tangent coordinates already absorbs much of the leading curvature
signal, so a degree-$2$-vs-degree-$2$ comparison cannot isolate the chart's contribution. The
two are not identical mechanisms --- a degree-$2$ tangent polynomial adds function-side
flexibility, whereas the second-order chart corrects the coordinate map itself --- but they
overlap heavily on the leading bias term, which is enough to make the old test uninformative.
The following isolates the curvature chart's value.

\subsection{E3.1 — Curvature-chart benefit under a degree-$1$, low-noise, curvature-stratified design}

\paragraph{Claim.} On genuinely curved manifolds, at \emph{degree~$1$} and adequate support, the
\code{second.order.svd} chart reduces Truth-RMSE in high-curvature regions relative to the
\code{pca} chart; and degree-$1$+\code{second.order.svd} approaches degree-$2$+\code{pca}
(confirming that degree-$2$ tangent polynomials make the chart correction \emph{largely redundant
in this regime} --- not that the two are identical mechanisms).

\paragraph{Failure mode addressed.} A confounded comparison (degree-$2$ both arms; noise-dominated;
whole-manifold averaging) that cannot reveal a curvature effect even if present.

\paragraph{DGP.} G3a, G3b (sphere cap), G3d (torus patch); curvature knob giving a low- and a
high-curvature setting per family; $n\in\{400,1600\}$; noise $\sigma\in\{0.01,0.03\}$ (low, so
bias is visible); $R=30$; plus the VALENCIA-derived real-geometry probe as an external check.

\paragraph{Configuration.} \emph{Matched} arms differing only in \code{local.chart.method}
$\in\{$\code{"pca"},\code{"second.order.svd"}$\}$. Fix degree $=1$, \code{chart.dim} $=$ true $d$,
the same kernel, and support $K\ge 3\cdot q_2$ where $q_2=d+\binom{d+1}{2}$ is the curvature
parameter count (for $d=2$, $q_2=5$, so $K\ge 20$; sweep $K\in\{25,40,60\}$). A third reference
arm: degree $2$ + \code{pca}, same $K$. Paired (Section~\ref{sec:paired}).

\paragraph{Statistic.} Local curvature magnitude $\kappa(x)$ is known analytically; bin eval
points into curvature deciles. Primary metric: Truth-RMSE within the top curvature decile,
$\rmse^{\mathrm{hi}}$. Secondary: whole-manifold Truth-RMSE; and
$\rmse(\text{deg1}+\text{2nd})-\rmse(\text{deg2}+\text{pca})$.

\paragraph{Acceptance / promotion criterion (predeclared).} Promote the curvature chart beyond
``experimental'' only if median over cases of $\Delta^{\mathrm{hi}}=\rmse^{\mathrm{hi}}_{\mathrm{pca}}-\rmse^{\mathrm{hi}}_{\mathrm{2nd}}$
is $>10\%$ relative with a sign-test $p<0.05$ across cases\,$\times$\,replicates, \emph{and}
runtime ratio is reported, \emph{and} the fallback rate (chart dim $=$ ambient, rank deficiency)
is $<10\%$ in the promoted regime. Separately, report whether degree-$1$+2nd $\approx$
degree-$2$+pca (redundancy confirmation).

\paragraph{Validity safeguards.} Degree fixed at $1$ in both compared arms (the central fix);
support sized for stable curvature estimation; low noise so bias dominates; curvature-stratified
metric so localized gains are not averaged away; matched seeds/support/kernel/dim; fallback
accounting mandatory; flat datasets excluded.

\subsection{E3.2 — Curvature-bias unit test (theory-anchored)}

\paragraph{Claim.} For an intrinsic-linear truth on a known-curvature manifold, degree-$1$+\code{pca}
exhibits the analytically predicted $O(\kappa h^2)$ bias, and degree-$1$+\code{second.order.svd}
removes its leading term.

\paragraph{Failure mode addressed.} A curvature-correction implementation that is plausible but
does not actually cancel the known curvature term (a correctness question independent of accuracy
gains).

\paragraph{DGP.} G3a with $f_{\mathrm{lin}}(u)=u_1$ (intrinsic-linear), noiseless; curvature
$\kappa=1/R$ swept over $R\in\{1,2,4,8\}$; support radius held fixed so $h$ is comparable across $R$.

\paragraph{Statistic.} At anchors in the high-curvature region, bias
$\beta^{(p,\mathrm{chart})}(x)=\hat f(x)-f(x)$; fit $|\bar\beta_{\mathrm{pca}}|$ vs.\ $\kappa h^2$
(expect linear through the origin); ratio $|\bar\beta_{\mathrm{2nd}}|/|\bar\beta_{\mathrm{pca}}|$.

\paragraph{Acceptance criterion.} $|\bar\beta_{\mathrm{pca}}|$ vs.\ $\kappa h^2$ has slope
significantly $>0$ and $R^2>0.9$ (the predicted scaling), and
$|\bar\beta_{\mathrm{2nd}}|/|\bar\beta_{\mathrm{pca}}|<0.3$ at the highest curvature (leading term
removed). If the chart does \emph{not} reduce the bias here, the correction is wrong regardless of
E3.1.

\paragraph{Validity safeguards.} Noiseless (pure bias); intrinsic-linear truth so degree-$1$ would
be exact \emph{but for} curvature, isolating the curvature term; $h$ held fixed across $R$ so the
$\kappa h^2$ scaling is identifiable.

\section{Tier 4 — Uncertainty quantification}

\subsection{E4.1 — Pointwise variance and confidence-band coverage}

\paragraph{Claim.} The linear-smoother variance $\Var(\hat y_i)=\sigma^2\sum_j S_{ij}^2$ yields
pointwise bands whose empirical coverage matches nominal in the interior; coverage degrades where
bias dominates (boundary/high curvature), as expected, and this is reported rather than masked.

\paragraph{Failure mode addressed.} Shipping point estimates with no uncertainty; or claiming
nominal coverage globally where bias guarantees undercoverage.

\paragraph{DGP.} G3a ($f_{\mathrm{smooth}}$), homoskedastic $\varepsilon\sim N(0,\sigma^2)$ with
$\sigma=0.1$ \emph{known}; $n=1200$; $R=500$ replicates (so coverage Monte-Carlo SE
$\approx\sqrt{0.95\cdot0.05/500}\approx 0.01$).

\paragraph{Configuration.} Fixed config (singleton grids), \code{chart.dim=2}, degree $1$. Build
$S$ (Section~\ref{sec:smoother}); band
$\hat y_i \pm z_{0.975}\,\hat\sigma\,\lVert S_{i\cdot}\rVert_2$ with $\hat\sigma^2=\mathrm{RSS}/(n-\tr S)$.

\paragraph{Statistic.} Per eval point, empirical coverage of $f(x_i)$ across replicates; stratify
into interior, boundary ($x$ within $h$ of the disk edge), and top-curvature-decile sets.

\paragraph{Acceptance criterion.} Interior average coverage $\in[0.93,0.97]$ (using known $\sigma$);
with the plug-in $\hat\sigma$, $\in[0.92,0.98]$. Boundary/high-curvature undercoverage is reported
with its magnitude; if a bias-corrected band is implemented, re-test that its interior$+$boundary
coverage both fall in range.

\paragraph{Validity safeguards.} Known $\sigma$ first (separates variance estimation from band
construction); honest stratification—do \emph{not} average interior and boundary coverage into one
headline number, since bias-driven boundary undercoverage is real and must be visible.

\section{Execution order and dependency map}

\begin{enumerate}[leftmargin=1.5em]
\item \textbf{Tier 0 first, as CI gates.} E0.2 (linear smoother) and E0.1 (reproduction) are the
keystones: E0.3a, E0.5, E4.1 reuse the $S$ machinery from E0.2, and any failure of E0.1/E0.2
invalidates everything downstream, so they run before any study.
\item \textbf{Tier 1 next}, because honest selection and bandwidth semantics underlie every
accuracy number reported in Tiers 3--4. E1.10 (nested/grouped CV) should be in place before any
real-data Truth-unknown claim.
\item \textbf{Tier 2} can proceed in parallel with Tier 1 (independent code paths).
\item \textbf{Tier 3} depends on E0.1 (reproduction on flat manifolds) and the matched-arm
protocol; E3.2 depends on E0.2's bias-extraction style.
\item \textbf{Tier 4} depends on E0.2 ($S$) and E0.5 (rate, for the bias/variance split).
\end{enumerate}

\paragraph{Per-experiment validity-threat summary.} Each experiment's ``Validity safeguards''
field is the controlling list. The recurring threats and their controls are: selection confounds
(use singleton grids for any $S$-based or bias-isolation test); bias floors (use the optimal
bandwidth schedule for rate tests); noise domination (use low $\sigma$ for bias-isolation tests);
averaging dilution (stratify by curvature/boundary); dependence leakage (grouped folds, disjoint
test clusters); and metric/deploy mismatch (score the deployed quantity). A result whose safeguards
were not met is reported as \emph{inconclusive}, never as evidence.

\end{document}
